\documentclass{ctexart}  % 移除 twocolumn

% 1. 基础设置与包
\usepackage[a4paper, margin=2cm, headheight=2.5cm, footskip=1cm, columnsep=1cm]{geometry}
\usepackage[dvipsnames]{xcolor}
\usepackage{graphicx}
\usepackage{amsmath}
\usepackage{fancyhdr}
\usepackage[absolute,overlay]{textpos}
\usepackage{enumitem}
\usepackage{dashrule}
\usepackage{framed}
\usepackage{multicol}  % 添加 multicol 包

% 2. 颜色定义
\definecolor{myblue}{rgb}{0.1, 0.3, 0.7}

% 3. 页面样式 (页眉与页脚)
\pagestyle{fancy}
\fancyhf{}

% --- 页眉 (章标题) ---
\fancyhead[L]{\songti\zihao{3} \textbf{第1章}}
\fancyhead[R]{\songti\zihao{3} \textbf{功和机械能}}

% --- 页脚 ---
% \fancyfoot[L]{\zihao{-5} 注: 2024表示2023\textasciitilde 2024学年, 2025表示2024\textasciitilde 2025学年.}
\fancyfoot[R]{\zihao{5} 努力不是为了给世界看的，而是为了去看世界。 \hspace{0.5em} \colorbox{gray!30}{\textbf{\thepage}}}

% 4. 自定义命令
% --- 题目来源 ---
\newcommand{\qsource}[1]{{\zihao{-5}\kaishu #1}\par}  % 使用 \kaishu 替代 \kai

% --- 多选标志 ---
\newcommand{\multichoice}{（多选）}

% --- 节标题 (居中, 蓝色, 虚线) ---
\newcommand{\sectiontitle}[1]{%
\end{multicols}  % 结束多栏
\par\noindent\filcenter
{\zihao{4} \textcolor{myblue}{#1}} \\
\vspace{1mm}
\noindent\filcenter
\textcolor{myblue}{\hdashrule{0.3\textwidth}{1pt}{2pt 2pt}}
\par\vspace{3mm}
\begin{multicols}{2}  % 重新开始多栏
}

% --- 板块标题 (刷基础) ---
\newcommand{\blocktitle}[1]{%
	\par\noindent
	{\zihao{4} \heiti \textcolor{myblue}{#1}}
	\par\vspace{2mm}
}

% --- 题型标题 ---
\newcommand{\qtype}[1]{%
	\par\noindent
	{\zihao{5} \heiti \textcolor{myblue}{#1}}
	\par\vspace{1mm}
}

% --- 题目列表环境 ---
\setlist[enumerate,1]{
	label=\textbf{\arabic*.},
	leftmargin=*,
	topsep=2mm,
	parsep=0pt,
	itemsep=1mm
}

% --- 选项列表环境 (A, B, C, D) ---
\setlist[enumerate,2]{
	label=\Alph*.,
	leftmargin=*,
	topsep=0.5ex,
	itemsep=0.2ex,
	parsep=0pt
}

% --------------------------------------------------
%  文档开始
% --------------------------------------------------
\begin{document}
	
	\begin{multicols}{2}  % 开始双栏
		
		\blocktitle{刷基础}
		
		% --- 题目区 (双栏) ---
		
		% --- 左栏 ---
		\qtype{题型1 对功的理解}
		
		\begin{enumerate}
			\item \qsource{[山东济南一中2025高一上月考] \multichoice}
			关于功，下列说法中正确的是（\quad）
			\begin{enumerate}
				\item 作用在物体上的力越大，做的功就越多
				\item 物体发生的位移越大，力做的功就越多
				\item 功是标量，但有正负，表示做功的效果
				\item 功是矢量，有大小和方向
			\end{enumerate}
			
			\item \qsource{李白《行路难》诗云："欲渡黄河冰塞川，将登太行雪满山。"}
			若某人（可视为质点）以恒定速率$v$登上高为$h$的山坡，在此过程中，人的（\quad）
			\begin{enumerate}
				\item 重力势能不变
				\item 动能不变
				\item 机械能守恒
				\item 合外力做功为 $mgh$
			\end{enumerate}
			
			\item \qsource{[福建福宁古五校教学联合体2025高一下联考]}
			如图所示，乒乓球以速度 $v$ 撞击球拍，与拍面成 $\theta$ 角，并以相同速率 $v$ 反弹，角度仍为 $\theta$。在碰撞瞬间，球拍对球的弹力做的功为（\quad）
			\begin{enumerate}
				\item $mv^2 \cos^2\theta$
				\item $mv^2 \sin^2\theta$
				\item $mv^2$
				\item 0
			\end{enumerate}

			\item \qsource{[福建福宁古五校教学联合体2025高一下联考]}
			如图所示，乒乓球以速度 $v$ 撞击球拍，与拍面成 $\theta$ 角，并以相同速率 $v$ 反弹，角度仍为 $\theta$。在碰撞瞬间，球拍对球的弹力做的功为（\quad）
			\begin{enumerate}
				\item $mv^2 \cos^2\theta$
				\item $mv^2 \sin^2\theta$
				\item $mv^2$
				\item 0
			\end{enumerate}

			\item \qsource{[福建福宁古五校教学联合体2025高一下联考]}
			如图所示，乒乓球以速度 $v$ 撞击球拍，与拍面成 $\theta$ 角，并以相同速率 $v$ 反弹，角度仍为 $\theta$。在碰撞瞬间，球拍对球的弹力做的功为（\quad）
			\begin{enumerate}
				\item $mv^2 \cos^2\theta$
				\item $mv^2 \sin^2\theta$
				\item $mv^2$
				\item 0
			\end{enumerate}
		\end{enumerate}
		
		% 强制换栏
		\columnbreak
		
		% --- 右栏 ---
		\qtype{题型2 恒力做功的计算}
		\begin{enumerate}[start=4]
			\item \qsource{[浙江杭州学军中学余杭学校2024高一下月考] \multichoice}
			如图所示，小孩坐在"海豚椅"上，在电梯中从1楼匀速上升到3楼。在此过程中（\quad）
			
			\begin{minipage}[t]{0.55\columnwidth}
				\begin{enumerate}
					\item 重力对小孩做正功
					\item 支持力对小孩做正功
					\item 合外力对小孩做功为零
					\item 小孩的机械能不变
				\end{enumerate}
			\end{minipage}
			
			\item \qsource{[山东淄博2025高一上期中]}
			如图，将四块质量均为 $m$、长度均为 $d$ 的相同砖块，从平铺状态（甲）变为叠放状态（乙），每块砖的重心在其中点。则人对砖块至少做的功为（\quad）
			\begin{enumerate}
				\item $2mgd$
				\item $4mgd$
				\item $6mgd$
				\item $8mgd$
			\end{enumerate}
			
			% \centering
			% \framebox[0.9\columnwidth][c]{\rule{0pt}{3cm}}
			
			% {\zihao{-5} (第5题图)}
			\par
			
			\item \qsource{[四川遂宁中学2025高一下期中] (教材变式)}
			一个质量为 $m=2\text{kg}$ 的物体，在 $F=10\text{N}$ 的水平拉力作用下，沿水平面移动了 $s=5\text{m}$。已知物体与地面间的动摩擦因数 $\mu=0.2$ ($g=10\text{m/s}^2$)。则拉力 $F$ 做的功为（\quad）
			\begin{enumerate}
				\item 10 J
				\item 20 J
				\item 30 J
				\item 50 J
			\end{enumerate}

			\item \qsource{[四川遂宁中学2025高一下期中] (教材变式)}
			一个质量为 $m=2\text{kg}$ 的物体，在 $F=10\text{N}$ 的水平拉力作用下，沿水平面移动了 $s=5\text{m}$。已知物体与地面间的动摩擦因数 $\mu=0.2$ ($g=10\text{m/s}^2$)。则拉力 $F$ 做的功为（\quad）
			\begin{enumerate}
				\item 10 J
				\item 20 J
				\item 30 J
				\item 50 J
			\end{enumerate}
		\end{enumerate}
		
	\end{multicols}  % 结束双栏
	
\end{document}